\section{Architekturentscheidungen}
\label{sec:adr}

Architekturentscheidungen (Architecture Decision Records, ADRs)
dokumentieren technische Entscheidungen mit Begründung und
bewerteten Alternativen.

\subsection{ADR-001: Kommunikationsprotokoll (UDP vs. RPC)}
\textbf{Betrifft}: Gesamte Systemkommunikation (Client, Dispatcher, Namensdienst, Worker)
\\\\
\noindent \textbf{Kontext}\\
\noindent Für die Kommunikation zwischen Client, Dispatcher, Namensdienst und Workern
wird ein Netzwerkprotokoll benötigt. Das System besteht aus mehreren
verteilten Komponenten, die unterschiedliche Nachrichten austauschen
(z. B. Task-Übermittlung, Worker-Registrierung, Heartbeats und Ergebnisrückgaben).
\\
Die Schnittstellen sollen klar definiert, typisiert und möglichst einfach
implementierbar sein. Außerdem soll die Kommunikation effizient erfolgen
und auch bei einer späteren Implementierung in anderen Programmiersprachen
weiterverwendbar bleiben.
\\\\
\textbf{Entscheidung}\\
\noindent \textbf{gRPC mit Protocol Buffers (RPC)} wurde als Kommunikationsprotokoll
für sämtliche Verbindungen zwischen den Komponenten gewählt.
\\
Alle Nachrichten werden als Protobuf-Nachrichten definiert und über
gRPC-Services ausgetauscht. Die Schnittstellendefinition erfolgt zentral
in der Datei \texttt{proto/taskgrid.proto}.

\subsubsection*{Bewertete Alternativen}

\begin{table}[H]
	\centering
	\begin{tabular}{m{2.5cm} m{6.75cm} m{6cm}}
		\toprule
		\textbf{Protokoll} & \textbf{Vorteile} & \textbf{Nachteile} \\
		\midrule
		\textbf{gRPC / RPC} (gewählt) &
		Strenge Typisierung, automatische Codegenerierung, effiziente binäre Serialisierung, sprachunabhängig &
		Höhere Komplexität als einfache Socket-Kommunikation \\
		\midrule
		UDP &
		Sehr geringer Overhead, schnelle Übertragung &
		Keine Garantie für Zustellung, Reihenfolge oder Duplikatfreiheit; Zuverlässigkeit müsste selbst implementiert werden \\
		\bottomrule
	\end{tabular}
	\caption{ADR-001: Bewertete Alternativen für das Kommunikationsprotokoll}
\end{table}

\subsubsection*{Begründung}

gRPC wurde gewählt, weil:

\begin{itemize}
	\item Das System viele strukturierte Nachrichten zwischen verteilten Komponenten austauscht.
	\item Die \texttt{.proto}-Datei als zentrale Schnittstellenbeschreibung dient und Inkonsistenzen zwischen Komponenten verhindert.
	\item Client, Dispatcher, Namensdienst und Worker unabhängig voneinander entwickelt werden können.
	\item Die automatische Generierung von Client- und Server-Code den Implementierungsaufwand reduziert.
	\item Zuverlässigkeit, Verbindungsmanagement und Fehlerbehandlung bereits durch gRPC bereitgestellt werden und nicht selbst implementiert werden müssen.
\end{itemize}

\subsubsection*{Konsequenzen}

Positiv: Einheitliche und typsichere Kommunikation im gesamten System,
geringer Serialisierungsaufwand, gute Erweiterbarkeit sowie einfache
Integration zusätzlicher Komponenten oder Programmiersprachen.
\\
Negativ: Abhängigkeit von gRPC- und Protobuf-Bibliotheken.
Die Kommunikation ist weniger transparent als JSON-basierte Protokolle
und Debugging auf Netzwerkebene kann aufwendiger sein.

\subsection{ADR-002: Heartbeat-Mechanismus}

\textbf{Betrifft}: Namensdienst, Worker-Management, Service Discovery
\\\\
\noindent\textbf{Kontext}\\
\noindent Der Dispatcher darf Tasks nur an erreichbare Worker senden.
Da Worker jederzeit abstürzen, beendet oder durch Netzwerkprobleme
unerreichbar werden können, muss der Namensdienst erkennen können,
welche Worker aktuell verfügbar sind.
\\
Eine einmalige Registrierung beim Start reicht hierfür nicht aus,
da der Namensdienst sonst nicht feststellen kann, ob ein Worker
zwischenzeitlich ausgefallen ist.
\\\\
\noindent \textbf{Entscheidung}\\
\noindent Es wird ein periodischer \textbf{Heartbeat-Mechanismus} eingesetzt.\\
Jeder Worker sendet nach erfolgreicher Registrierung in festen
Zeitabständen einen Heartbeat an den Namensdienst.
Bleibt ein Heartbeat länger als die konfigurierte Timeout-Schwelle aus,
wird der Worker automatisch als inaktiv markiert und aus der Liste
verfügbarer Worker entfernt.

\begin{itemize}
	\item Heartbeat-Intervall: \texttt{HEARTBEAT\_INTERVAL\_SECONDS}
	\item Timeout-Schwelle: \texttt{HEARTBEAT\_TIMEOUT\_SECONDS}
	\item Zustandsänderung: \texttt{ACTIVE} $\rightarrow$ \texttt{INACTIVE}
	\item Inaktive Worker werden von \texttt{LookupWorker()} nicht mehr zurückgegeben.
\end{itemize}

\subsubsection*{Bewertete Alternativen}

\begin{table}[H]
	\centering
	\begin{tabular}{m{3.25cm} m{6.25cm} m{6.5cm}}
		\toprule
		\textbf{Verfahren} & \textbf{Vorteile} & \textbf{Nachteile} \\
		\midrule
		\textbf{Heartbeat} &
		Einfach implementierbar, schnelle Fehlererkennung, keine aktive Überwachung durch Dispatcher erforderlich &
		Erzeugt regelmäßigen Netzwerkverkehr \\
		\midrule
		Periodisches Polling durch Namensdienst &
		Kontrolle liegt zentral beim Namensdienst &
		Zusätzliche Verbindungen; Namensdienst muss alle Worker aktiv kontaktieren \\
		\midrule
		Statische Registrierung &
		Minimaler Kommunikationsaufwand &
		Ausfälle werden nicht erkannt; veraltete Worker-Einträge bleiben bestehen \\
		\bottomrule
	\end{tabular}
	\caption{ADR-005: Bewertete Alternativen für Worker-Überwachung}
\end{table}

\subsubsection*{Begründung}

Der Heartbeat-Mechanismus wurde gewählt, weil:

\begin{itemize}
	\item Worker dynamisch gestartet und beendet werden können.
	\item Der Namensdienst jederzeit eine aktuelle Liste verfügbarer Worker benötigt.
	\item Die Implementierung einfach und unabhängig vom Dispatcher ist.
	\item Ausfälle automatisch erkannt werden, ohne dass andere Komponenten eingreifen müssen.
	\item Das Verfahren auch bei mehreren Worker-Instanzen gut skaliert.
\end{itemize}

\subsubsection*{Konsequenzen}

Positiv: Ausgefallene Worker werden automatisch erkannt und aus der
Service-Discovery entfernt. Der Dispatcher erhält nur noch erreichbare
Worker und kann Tasks zuverlässig verteilen.
\\
Negativ: Zwischen Ausfall und Erkennung vergeht eine gewisse Zeit,
die durch Heartbeat-Intervall und Timeout-Schwelle bestimmt wird.
Kürzere Intervalle verbessern die Reaktionszeit, erhöhen jedoch die
Netzwerklast und die Anzahl der Heartbeat-Nachrichten.


\subsection{ADR-003: Worker-Auswahlstrategie - Round-Robin}

\textbf{Status}: Akzeptiert \\
\textbf{Betrifft}: \texttt{src/dispatcher/worker\_selector.py} (Issue \#14, \#15)

\subsubsection*{Kontext}

Wenn der Namensdienst mehrere Worker für einen Task-Typ zurückgibt,
muss der Dispatcher einen davon auswählen.
Die Aufgabenstellung verbietet, dauerhaft denselben Worker zu wählen.
Es wird eine Strategie benötigt die einfach, korrekt und skalierbar ist.

\subsubsection*{Entscheidung}

\textbf{Round-Robin}: Pro Task-Typ wird ein Zähler geführt.
Bei $n$ verfügbaren Workern wird Worker mit Index $\text{Zähler} \bmod n$
gewählt, dann der Zähler inkrementiert.

\begin{lstlisting}[language=Python,caption={Round-Robin-Implementierung}]
with self._lock:
    idx = self._counters[task_type] % len(workers)
    self._counters[task_type] += 1
return workers[idx]
\end{lstlisting}

\subsubsection*{Bewertete Alternativen}

\begin{table}[H]
\centering
\begin{tabular}{lp{5cm}p{5cm}}
\toprule
\textbf{Strategie} & \textbf{Vorteile} & \textbf{Nachteile} \\
\midrule
\textbf{Round-Robin} (gewählt) &
  Einfach, deterministisch, gleichmäßige Last, kein globaler Zustand außer Index &
  Ignoriert aktuelle Auslastung der Worker \\
\midrule
Random-Auswahl &
  Keine Synchronisation nötig &
  Kann zufällig denselben Worker mehrfach wählen; nicht deterministisch für Tests \\
\midrule
Least-Load &
  Wählt immer den am wenigsten ausgelasteten Worker &
  Benötigt aktuelle Auslastungsdaten; Namensdienst muss diese liefern;
  Race Condition wenn mehrere Dispatcher-Threads gleichzeitig wählen \\
\midrule
Consistent Hashing &
  Gleiche Tasks gehen immer zum gleichen Worker (Cache-freundlich) &
  Aufwändig; unnötig für zustandslose Worker \\
\bottomrule
\end{tabular}
\caption{ADR-003: Bewertete Alternativen für Worker-Auswahl}
\end{table}

\subsubsection*{Begründung}

Round-Robin wurde gewählt, weil:
\begin{itemize}
  \item Die Aufgabenstellung explizit eine Verteilungsstrategie fordert
        (kein immer-derselbe-Worker).
  \item Die Implementierung minimal ist (ein Zähler-Dictionary, ein Lock).
  \item Das Ergebnis bei Tests deterministisch und nachvollziehbar ist.
  \item Die Worker in diesem System zustandslos sind -- Least-Load
        brächte keinen nennenswerten Vorteil, da jeder Task unabhängig
        verarbeitet wird.
\end{itemize}

\subsubsection*{Konsequenzen}

Positiv: Einfache Implementierung, keine zyklische Abhängigkeit zwischen
Dispatcher und Auslastungsdaten der Worker.
Negativ: Bei stark unterschiedlicher Task-Dauer kann ein Worker überlastet
werden während andere idle sind. In diesem System ist das akzeptabel,
da die Aufgaben (reverse, sum, hash, upper, wait) ähnliche Laufzeiten haben.

\subsection{ADR-004: Timeout-Behandlung und Retry-Strategie}

\textbf{Status}: Akzeptiert \\
\textbf{Betrifft}: \texttt{src/dispatcher/dispatch\_loop.py} (Issue \#18)

\subsubsection*{Kontext}

Ein Task kann nach dem Versenden an einen Worker nicht zurückkommen, wenn:
\begin{itemize}
  \item Der Worker-Prozess abstürzt oder hängt.
  \item Der Netzwerkpfad zum Worker unterbrochen ist.
  \item Die Verarbeitungszeit das konfigurierte Timeout überschreitet.
\end{itemize}

Der Dispatcher muss erkennen, wenn kein Ergebnis eintrifft, und
angemessen reagieren. Die Aufgabenstellung fordert eine Dokumentation
der gewählten Strategie.

\subsubsection*{Entscheidung}

\textbf{Timer-basiertes Timeout mit konfigurierbarem Re-Enqueue + Max-Retries}:

\begin{enumerate}
  \item Beim Versenden eines Tasks an einen Worker startet ein
        \texttt{threading.Timer} mit \texttt{DISPATCH\_TIMEOUT\_SECONDS}.
  \item Feuert der Timer (kein Ergebnis eingegangen):
        DISPATCHED $\rightarrow$ TIMEOUT $\rightarrow$ RETRYING,
        \texttt{retry\_count++}.
  \item Falls \texttt{retry\_count} $<$ \texttt{MAX\_RETRIES}:
        Task wird erneut in die Queue eingereiht (RETRYING).
  \item Falls \texttt{retry\_count} $\geq$ \texttt{MAX\_RETRIES}:
        RETRYING $\rightarrow$ FAILED (endgültig).
  \item Trifft ein Ergebnis vor dem Timeout ein: Timer wird sofort
        abgebrochen.
  \item Verspätetes Feuern auf terminalen Task: wird ignoriert.
\end{enumerate}

\subsubsection*{Bewertete Alternativen}

\begin{table}[H]
\centering
\begin{tabular}{m{4.5cm} m{5.2cm} m{6.25cm}}
\toprule
\textbf{Strategie} & \textbf{Vorteile} & \textbf{Nachteile} \\
\midrule
\textbf{Timer + Re-Enqueue} (gewählt) &
  Saubere Zustandsübergänge; konfigurierbares Max-Retry;
  Timer-Abbruch bei rechtzeitigem Ergebnis &
  Mehrere Threads (Timer-Threads); Race Condition möglich
  wenn Ergebnis und Timeout gleichzeitig eintreffen
  (durch is\_terminal()-Check behoben) \\
\midrule
Sofortiges FAILED &
  Einfach; kein Retry-Zustand nötig &
  Kein Retry bei transientem Fehler; schlechte Robustheit \\
\midrule
Weitersenden an anderen Worker (ohne Zustandsübergang) &
  Schneller als Re-Enqueue &
  Umgeht die Zustandsmaschine; schwerer testbar;
  Timeout-Zustand fehlt \\
\midrule
Polling statt Timer &
  Kein separater Thread pro Task &
  Hohe CPU-Last durch ständiges Polling;
  Reaktionszeit von Poll-Intervall abhängig \\
\bottomrule
\end{tabular}
\caption{ADR-005: Bewertete Alternativen für Timeout-Behandlung}
\end{table}

\subsubsection*{Begründung}

Die Timer-basierte Strategie wurde gewählt, weil:
\begin{itemize}
  \item Sie die vollständige Zustandsmaschine (\texttt{TIMEOUT},
        \texttt{RETRYING}) durchläuft -- alle 8 geforderten Zustände
        sind damit erreichbar.
  \item \texttt{threading.Timer} ist Teil der Python-Standardbibliothek,
        kein externes Framework nötig.
  \item Die Konfigurierbarkeit (\texttt{DISPATCH\_TIMEOUT\_SECONDS},
        \texttt{MAX\_RETRIES}) ermöglicht Anpassung ohne Code-Änderung.
  \item Die \texttt{is\_terminal()}-Prüfung im Timer-Callback löst die
        Race Condition zwischen gleichzeitig eintreffendem Ergebnis
        und Timeout-Feuern.
\end{itemize}

\subsubsection*{Konfiguration}

\begin{table}[H]
\centering
\begin{tabular}{lll}
\toprule
\textbf{Variable} & \textbf{Default} & \textbf{Bedeutung} \\
\midrule
\texttt{DISPATCH\_TIMEOUT\_SECONDS}        & 30 & Sekunden bis Timeout-Feuern \\
\texttt{MAX\_RETRIES}                      & 3  & Max. Retry-Versuche vor FAILED \\
\texttt{DISPATCH\_NO\_WORKER\_RETRY\_SECS} & 2  & Wartezeit wenn kein Worker verfügbar \\
\bottomrule
\end{tabular}
\caption{ADR-005: Timeout-Konfiguration}
\end{table}

\subsubsection*{Konsequenzen}

Positiv: Robustheit gegenüber ausgefallenen Workern.
Ein transienter Fehler (Worker kurz überlastet, Netzwerk-Hickup)
führt nicht sofort zu FAILED.
Das Log zeigt jeden Retry mit \texttt{retry\_count},
was die Fehlerdiagnose erleichtert.

Negativ: Im schlechtesten Fall wartet ein Task
$\texttt{MAX\_RETRIES} \times \texttt{DISPATCH\_TIMEOUT\_SECONDS}
= 3 \times 30 = 90$ Sekunden, bevor er als FAILED markiert wird.
Für Anwendungen mit harten Latenzanforderungen müsste der Timeout
entsprechend kleiner gewählt werden.

\subsection{ADR-006: Logging-Format und Propagation}

\textbf{Status}: Akzeptiert \\
\textbf{Betrifft}: \texttt{src/common/logger.py} (Issue \#19)

\subsubsection*{Kontext}

Das Logging-Modul wird von allen Dispatcher-Komponenten genutzt.
Python-Logger propagieren Nachrichten standardmäßig aufwärts in der
Logger-Hierarchie. Bei mehreren Komponenten mit eigenen Loggern
(\texttt{dispatcher}, \texttt{dispatcher.dispatch\_loop} etc.) führt dies
zu doppelten oder mehrfachen Log-Ausgaben.

\subsubsection*{Entscheidung}

\texttt{logger.propagate = False} wird in \texttt{get\_logger()} gesetzt,
sobald Handler hinzugefügt wurden.

\subsubsection*{Begründung}

Ohne diese Einstellung würde ein Log-Eintrag in
\texttt{dispatcher.dispatch\_loop} auch durch den \texttt{dispatcher}-Logger
und den Root-Logger verarbeitet -- je nachdem wie viele Handler
registriert sind, führt dies zu 2--4 identischen Ausgaben pro Event.
Das Testprotokoll würde unlesbar.

\subsubsection*{Konsequenzen}

Jede Komponente (\texttt{dispatcher}, \texttt{dispatcher.dispatch\_loop},
\texttt{dispatcher.namensdienst\_client} etc.) schreibt ihre Logs
ausschließlich in ihre eigene Datei und auf stdout.
Keine Überschneidungen, keine doppelten Einträge.
